% Define custom counters
\newcounter{sectiontoc}
\newcounter{subsectiontoc}
\newcounter{subsubsectiontoc}
\newcounter{paragraphtoc}

% appendix utils
% Define a flag for appendix
\newbool{inappendix}
\setbool{inappendix}{false}
% Redefine \appendix to set the flag
\let\oldappendix\appendix
\renewcommand{\appendix}{
  \setbool{inappendix}{true}
  \oldappendix
}

% title and toc customization

% First level: 第\chinese{section}章
% Second level: 一、
% Third level: 1.
% Fourth level: (1)
% Increment counters in titlecontents with condition
\titlecontents{section}[0pt]{\normalfont\normalsize}{
  \ifnum\value{sectiontoc}>4 % Condition: skip numbering if sectiontoc is greater than 2
    \stepcounter{sectiontoc} % Increment counter without displaying\
    \hspace{1em} % Adjust spacing for unnumbered section
  \else
    \stepcounter{sectiontoc}\hspace{1em}第\chinese{sectiontoc}章\ 
  \fi
  \setcounter{subsectiontoc}{0}\
  \setcounter{subsubsectiontoc}{0}\
  \setcounter{paragraphtoc}{0}\
}{\hspace{1.5em}}{\titlerule*[0.5pc]{.}\contentspage}

\titlecontents{subsection}[2em]{\normalfont\normalsize\sffamily}{
  \stepcounter{subsectiontoc}\chinese{subsectiontoc}、\
  \setcounter{subsubsectiontoc}{0}\
  \setcounter{paragraphtoc}{0}\
}{1em}{\titlerule*[0.5pc]{.}\contentspage}

\titlecontents{subsubsection}[2.5em]{\normalfont\normalsize\sffamily}{
  \stepcounter{subsubsectiontoc}\arabic{subsubsectiontoc}.\
  \setcounter{paragraphtoc}{0}\
}{1em}{\titlerule*[0.5pc]{.}\contentspage}

\titlecontents{paragraph}[3em]{\normalfont\normalsize\sffamily}{
  \stepcounter{paragraphtoc}(\arabic{paragraphtoc})\
}{1em}{\titlerule*[0.5pc]{.}\contentspage}

% First level: 小三、黑体、居中、1.5倍行间距
\titleformat{\section}
  {\centering\normalfont\sffamily\zihao{-3}}
  {\ifbool{inappendix}{}{第\chinese{section}章\hspace{1em}}}
  {0pt}{}
%\titlespacing{\section}{0pt}{*2}{*1.5}
\titlespacing{\section}{0pt}{1\baselineskip}{1\baselineskip}
% Second level: 四号、黑体、顶头、1.5倍行间距
\titleformat{\subsection}{\normalfont\sffamily\zihao{4}}{\chinese{subsection}、}{1em}{}
%\titlespacing{\subsection}{0pt}{*2}{*1.5}
%\titlespacing{\subsection}{0pt}{1.5\baselineskip}{1.5\baselineskip}
% Third and Fourth levels: 小四、宋体、1.5倍行距
\titleformat{\subsubsection}{\normalfont\sffamily\zihao{-4}}{\arabic{subsubsection}.}{1em}{}
%\titlespacing{\subsubsection}{0pt}{*2}{*1.5}
%\titlespacing{\subsubsection}{0pt}{1.5\baselineskip}{1.5\baselineskip}
\titleformat{\paragraph}[runin]{\normalfont\sffamily\zihao{-4}}{(\arabic{paragraph})}{1em}{}
%\titlespacing{\paragraph}{0pt}{*2}{*1.5}
%\titlespacing{\paragraph}{0pt}{1.5\baselineskip}{1.5\baselineskip}
% Set depth for numbering and TOC
\setcounter{secnumdepth}{4}
\setcounter{tocdepth}{3}

% 设置正文格式
%\doublespacing
%\onehalfspacing
%\normalsize
%\renewcommand{\arraystretch}{1.5}
%\renewcommand{\baselinestretch}{1.5}
\renewcommand{\arraystretch}{1.5}
\renewcommand{\baselinestretch}{1.5}

\setlength{\parindent}{2em} % 缩进
%\renewcommand{\rmdefault}{song} % 宋体

%\setwordspacing
% Custom formats for the table of contents
% Remove numbering in the table of contents
% \renewcommand{\cftsecpresnum}{}
% \renewcommand{\cftsecaftersnum}{}
% \renewcommand{\cftsubsecpresnum}{}
% \renewcommand{\cftsubsecaftersnum}{}

\ExplSyntaxOn
% This macro based on https://tex.stackexchange.com/a/25704/
\NewDocumentCommand {\hfillspace}
  { > { \SplitList { ~ } } m }
  { \tl_map_inline:nn {#1} { ##1\hfill } }
\ExplSyntaxOff
\newcounter{myline}
\AtBeginEnvironment{tabular}{\setcounter{myline}{0}}
%\begin{tabular}{>{\stepcounter{myline}\themyline\quad\collectcell\hfillspace}p{.9\columnwidth}<{\endcollectcell}}
%font
% 字体已在 00-depends.tex 中统一设置（含自动检测与回退），此处不再重复，以避免未安装字体时报错。
% 如需局部改变标题字体，可使用 \TitleCJKFont 命令包装。

%\setCJKsansfont{NotoSansMNerdFont-Regular} % 设置黑体为无衬线字体
%\setCJKmainfont{NotoSerifNerdFont-Regular} % 设置宋体为主字体
%\renewcommand{\cfttoctitlefont}{\hfill\Huge\bfseries}
%\renewcommand{\cftaftertoctitle}{\hfill}
\renewcommand{\contentsname}{ \centering{目\hspace{1em}录}}
%\renewcommand{\contentsname}{\centering Table of Contents}
\newcommand{\fixedunderline}[2][6cm]{\underline{\makebox[#1][c]{#2}}}

% code block
% example:
% \lstinputlisting{/path/to/your/codefile.m}
\lstset{
  language=Matlab, % or the language of your code
  basicstyle=\ttfamily, % Set the font style
  numbers=left, % Add line numbers on the left
  numberstyle=\tiny, % Style for line numbers
  stepnumber=1, % Line number step
  numbersep=5pt, % Distance from code
  showspaces=false, % Don't show spaces
  showstringspaces=false, % Don't show spaces in strings
  showtabs=false, % Don't show tabs
  frame=single, % Frame around code
  tabsize=2, % Size of tabs
  breaklines=true, % Enable line breaking
  breakatwhitespace=true, % Break at whitespace
  captionpos=b, % Position of caption
  keepspaces=true, % Keep spaces in text
  columns=flexible, % Adjust column width
}

